\chapter{距離空間・連続・コンパクト性}
\label{ch:found-metric}

\begin{definition}{距離空間}{sem-metric-space}
\blockmeta{space=metric}
  集合 $\Omega$ と関数 $d \colon \Omega \times \Omega \to [0, \infty)$ の組
  $(\Omega, d)$ が\textbf{距離空間}であるとは，
  次の3条件がすべての $x, y, z \in \Omega$ に対して成り立つことをいう：
  \begin{enumerate}
    \item[(i)] 同一性：$d(x, y) = 0 \iff x = y$，
    \item[(ii)] 対称性：$d(x, y) = d(y, x)$，
    \item[(iii)] 三角不等式：$d(x, z) \leq d(x, y) + d(y, z)$．
  \end{enumerate}
  関数 $d$ を $\Omega$ 上の\textbf{距離関数}という．
\end{definition}

\begin{definition}{ノルム空間}{sem-normed-space}
\blockmeta{space=normed}
  実ベクトル空間 $V$ と関数 $\|\cdot\| \colon V \to [0, \infty)$ の組
  $(V, \|\cdot\|)$ が\textbf{ノルム空間}であるとは，
  次の3条件がすべての $x, y \in V$ と $\lambda \in \R$ に対して
  成り立つことをいう：
  \begin{enumerate}
    \item[(i)] 正定値性：$\|x\| = 0 \iff x = 0$，
    \item[(ii)] 斉次性：$\|\lambda x\| = |\lambda|\,\|x\|$，
    \item[(iii)] 三角不等式：$\|x + y\| \leq \|x\| + \|y\|$．
  \end{enumerate}
  関数 $\|\cdot\|$ を $V$ 上の\textbf{ノルム}という．
  ノルム空間は $d(x, y) \defeq \|x - y\|$ により
  距離空間（定義~\ref{def:sem-metric-space}）となる．
  $V$ が有限次元のとき，$(V, \|\cdot\|)$ を
  \textbf{有限次元ノルム空間}という．
  例えば $\R^n$ や $\R^{n \times m}$ はユークリッドノルム（フロベニウスノルム）により
  有限次元ノルム空間となる．
\end{definition}

\begin{definition}{開球}{sem-open-ball}
  距離空間（定義~\ref{def:sem-metric-space}）$(\Omega, d)$ 上の点 $x \in \Omega$ と半径 $r \in \Rpos$ に対して，
  中心 $x$，半径 $r$ の\textbf{開球}を
  \[
    B_r(x) \defeq \{y \in \Omega \mid d(x, y) < r\}
  \]
  で定める．
\end{definition}

距離空間 $(\Omega, d)$ における開集合（定義~\ref{def:sem-topological-space}）は，
開球 $B_r(x)$ を用いて次のように特徴付けられる：
$U \subset \Omega$ が開集合であるとは，任意の $x \in U$ に対し
$B_r(x) \subset U$ をみたす $r \in \Rpos$ が存在することをいう．

\begin{figure}[H]
\centering
\begin{tikzpicture}[scale=1.0]
  %% --- 左：開集合 U（境界は U に含まれない） ---
  \begin{scope}
    \node[above, font=\small\bfseries] at (1.7, 3.0) {開集合の例};
    %% U の領域（境界を点線で示し，「境界は含まない」）
    \draw[blue, dashed, thick, fill=blue!8]
      plot[smooth cycle, tension=0.7]
      coordinates {(0.2, 1.0) (0.7, 2.4) (2.1, 2.7) (3.2, 2.0)
                   (3.0, 0.8) (1.8, 0.2) (0.6, 0.4)};
    \node[blue, font=\small] at (1.9, 1.5) {$U$};

    %% 内点 x_1 とその開球
    \fill[red] (1.2, 1.7) circle (1.5pt);
    \node[red, font=\small, above right] at (1.2, 1.7) {$x_1$};
    \draw[red, thick] (1.2, 1.7) circle (0.35);
    \node[red, font=\footnotesize] at (1.2, 1.05) {$B_{r_1}(x_1) \subset U$};

    %% 内点 x_2 とその開球（境界に近い）
    \fill[red] (2.6, 2.0) circle (1.5pt);
    \node[red, font=\small, above] at (2.6, 2.05) {$x_2$};
    \draw[red, thick] (2.6, 2.0) circle (0.18);

    \node[font=\footnotesize, anchor=north] at (1.7, -0.1)
      {任意の点 $x \in U$ で $B_r(x) \subset U$ なる $r$ がとれる};
  \end{scope}

  %% --- 右：開集合でない例（境界点を含む） ---
  \begin{scope}[xshift=6.0cm]
    \node[above, font=\small\bfseries] at (1.7, 3.0) {開集合でない例};
    %% F の領域（境界を実線で示し，「境界も含む」）
    \draw[red!70!black, thick, fill=red!8]
      plot[smooth cycle, tension=0.7]
      coordinates {(0.2, 1.0) (0.7, 2.4) (2.1, 2.7) (3.2, 2.0)
                   (3.0, 0.8) (1.8, 0.2) (0.6, 0.4)};
    \node[red!70!black, font=\small] at (1.9, 1.5) {$F$};

    %% 境界上の点 x_0
    \fill[red] (3.18, 1.7) circle (1.5pt);
    \node[red, font=\small, right] at (3.22, 1.7) {$x_0$};
    %% B_r(x_0) が F からはみ出す様子
    \draw[red, thick, dashed] (3.18, 1.7) circle (0.4);
    \node[red, font=\footnotesize] at (3.6, 1.05) {$B_r(x_0) \not\subset F$};

    \node[font=\footnotesize, anchor=north] at (1.7, -0.1)
      {境界点 $x_0$ で $B_r(x_0)$ が外にはみ出す};
  \end{scope}
\end{tikzpicture}
\caption{開集合の概念図．左：$U$ は開集合で，任意の点 $x \in U$ について
十分小さい開球 $B_r(x)$ が $U$ に含まれる（境界は $U$ に含まれない；点線で表示）．
右：$F$ は境界を含むため開集合でなく，境界点 $x_0$ ではどんなに $r$ を
小さくしても $B_r(x_0) \not\subset F$．
$\Omega$ 上の開集合全体が位相 $\mathcal{O}_d$ をなす．}
\label{fig:sem-open-set}
\end{figure}

\begin{definition}{収束列}{sem-convergence}
  距離空間 $(\Omega, d)$ における列 $(x_n)_{n \geq 1} \subset \Omega$ が
  点 $x \in \Omega$ に\textbf{収束}するとは，
  任意の $\varepsilon > 0$ に対してある $N \in \N$ が存在して
  $n \geq N \Longrightarrow d(x_n, x) < \varepsilon$ となることをいう．
  このとき $\lim_{n \to \infty} x_n = x$ または $x_n \to x$ と書く．
\end{definition}

\begin{definition}{完備性}{sem-complete-metric}
\blockmeta{space=complete_metric}
  距離空間 $(\Omega, d)$ の列 $(x_n)_{n \geq 1} \subset \Omega$ が
  \textbf{Cauchy 列}であるとは，
  任意の $\varepsilon > 0$ に対してある $N \in \N$ が存在して
  $n, m \geq N \Longrightarrow d(x_n, x_m) < \varepsilon$ となることをいう．
  $(\Omega, d)$ が\textbf{完備}であるとは，
  $\Omega$ における任意の Cauchy 列が $\Omega$ の点に収束する
  （定義~\ref{def:sem-convergence}）ことをいう．すなわち
  \[
    \forall (x_n)_{n \geq 1} \subset \Omega \text{ Cauchy 列},\;
    \exists x \in \Omega,\;
    x_n \to x.
  \]
  収束の定義に展開すれば
  \[
    \forall (x_n)_{n \geq 1} \subset \Omega \text{ Cauchy 列},\;
    \exists x \in \Omega,\;
    \forall \varepsilon > 0,\;
    \exists N \in \N,\;
    n \geq N \Longrightarrow d(x_n, x) < \varepsilon.
  \]
\end{definition}

\begin{figure}[H]
\centering
\begin{tikzpicture}[scale=1.0]
  %% --- 左：完備な空間 Ω = [0, 4] ---
  \begin{scope}
    \node[above, font=\small\bfseries] at (2.2, 1.7) {完備：$\Omega = [0, 4]$};
    \draw[thick] (0, 1.0) -- (4.2, 1.0);
    \fill[black] (0, 1.0) circle (1.5pt);
    \fill[black] (4.2, 1.0) circle (1.5pt);

    %% Cauchy sequence x_k = 1/k * 3 (decreasing toward 0)
    \foreach \k/\xpos in {1/3.0, 2/1.5, 3/1.0, 4/0.75, 5/0.6, 6/0.5} {
      \fill[blue] (\xpos, 1.0) circle (1.5pt);
    }
    \node[blue, font=\scriptsize, above] at (3.0, 1.1) {$x_1$};
    \node[blue, font=\scriptsize, above] at (1.5, 1.1) {$x_2$};
    \node[blue, font=\scriptsize, above] at (1.0, 1.1) {$x_3$};
    \node[blue, font=\footnotesize, above] at (0.5, 1.3) {$\cdots$};

    %% limit point at 0 ∈ Ω
    \fill[red] (0, 1.0) circle (2.5pt);
    \node[red, font=\small, below] at (0, 0.85) {$x^*$};
    \node[red, font=\footnotesize, below] at (0, 0.5) {$\in \Omega$};

    %% convergence arrow
    \draw[->, gray, thick]
      (1.0, 0.45) .. controls (0.5, 0.5) .. (0.05, 0.95);
    \node[font=\footnotesize, gray, anchor=west] at (1.1, 0.4) {収束先};
  \end{scope}

  %% --- 右：完備でない $\Omega' = (0, 4]$ ---
  \begin{scope}[xshift=5.5cm]
    \node[above, font=\small\bfseries] at (2.2, 1.7) {完備でない：$\Omega' = (0, 4]$};
    \draw[thick] (0.05, 1.0) -- (4.2, 1.0);
    %% 0 を除く: open endpoint
    \draw[thick, fill=white] (0, 1.0) circle (2.0pt);
    \fill[black] (4.2, 1.0) circle (1.5pt);

    %% same Cauchy 列
    \foreach \xpos in {3.0, 1.5, 1.0, 0.75, 0.6, 0.5} {
      \fill[blue] (\xpos, 1.0) circle (1.5pt);
    }

    %% 極限が Ω' の外
    \draw[red, very thick] (-0.18, 0.83) -- (0.18, 1.17);
    \draw[red, very thick] (-0.18, 1.17) -- (0.18, 0.83);
    \node[red, font=\small, below] at (0, 0.7) {$x^*$};
    \node[red, font=\footnotesize, below] at (0, 0.4) {$\notin \Omega'$};
  \end{scope}
\end{tikzpicture}
\caption{完備性の概念．Cauchy 列 $(x_k)$ は項間距離が無限小に縮まる
（$\forall \varepsilon > 0,\; \exists N,\; k, l \geq N \Rightarrow d(x_k, x_l) < \varepsilon$）．
左：完備な $\Omega$ では極限 $x^*$ が $\Omega$ 内に存在し $x_k \to x^*$．
右：$\Omega' = (0, 4]$（端点 $0$ を除外）では同じ Cauchy 列の極限 $0$ が
$\Omega'$ の外に逃げ，完備でない．}
\label{fig:sem-completeness}
\end{figure}

\begin{definition}{Polish 空間}{sem-polish}
\blockmeta{space=polish}
  完備（定義~\ref{def:sem-complete-metric}）かつ
  可分（定義~\ref{def:sem-separable}）な
  距離空間（定義~\ref{def:sem-metric-space}）を\textbf{Polish 空間}という．
  例：$\R^d$ はユークリッド距離に関して Polish 空間である
  （$\Q^d$ が可算稠密部分集合を与える）．
\end{definition}

\begin{figure}[H]
\centering
\begin{tikzpicture}[scale=1.0]
  %% Ω = R
  \draw[thick, ->] (-0.5, 0) -- (8.0, 0);
  \node[right, font=\small] at (8.0, 0) {$\R$};

  %% Q = countable dense subset, labeled
  %% spread Q points; some near limit x* = π
  \foreach \k/\xpos in {1/0.5, 2/1.5, 3/4.0, 4/6.5, 5/7.3} {
    \fill[blue] (\xpos, 0) circle (1.5pt);
    \node[blue, font=\scriptsize, below] at (\xpos, -0.15) {$q_{\k}$};
  }
  %% Cauchy subsequence (q_{k_n}) → x* (irrational point)
  \def\xstar{3.14}
  \foreach \i/\xpos/\lab in {6/2.5/$q_{k_1}$, 7/3.0/$q_{k_2}$, 8/3.13/$q_{k_3}$} {
    \fill[blue] (\xpos, 0) circle (1.5pt);
    \node[blue, font=\scriptsize, above] at (\xpos, 0.1) {\lab};
  }
  %% extra Q points (background density)
  \foreach \xpos in {1.0, 1.9, 4.5, 5.0, 5.5, 6.0, 6.9, 7.6} {
    \fill[blue!50] (\xpos, 0) circle (1pt);
  }

  %% x* (limit, irrational) in R
  \fill[red] (\xstar, 0) circle (2.5pt);
  \node[red, font=\small, above] at (\xstar, 0.3) {$x^*$};
  \node[red, font=\scriptsize] at (\xstar, 0.6) {（無理数）};

  %% arrow from Cauchy seq to x*
  \draw[->, blue!50!green, thick]
    (2.7, 1.3) .. controls (3.0, 1.0) .. (\xstar, 0.4);
  \node[blue!50!green, font=\footnotesize, anchor=south west] at (1.3, 1.25)
    {Cauchy 列 $(q_{k_n}) \to x^*$};

  %% labels
  \node[font=\footnotesize, anchor=west, blue] at (-0.3, -1.0)
    {\textbf{可分性}：可算稠密な $\Q = \{q_1, q_2, \ldots\} \subset \R$};
  \node[font=\footnotesize, anchor=west, red!70!black] at (-0.3, -1.4)
    {\textbf{完備性}：Cauchy 列の極限 $x^*$ が $\R$ 内に存在
      （無理数も $\R$ には含まれる）};
\end{tikzpicture}
\caption{Polish 空間（完備可分距離空間）のイメージ．代表例 $\R$ 上で：
$\Q$ は可算かつ稠密（可分性）．
$\Q$ 内の Cauchy 列 $(q_{k_n})$ の極限は無理数になりうるが，
それも $\R$ 内に存在する（完備性）．
$\Q$ 自身は完備でない（極限を $\Q$ に拾えない）が，
$\R$ はその「完備化」として機能する．}
\label{fig:sem-polish}
\end{figure}

\begin{definition}{連続写像}{sem-continuous}
  距離空間 $(\Omega_1, d_1)$，$(\Omega_2, d_2)$ の間の写像
  $f \colon \Omega_1 \to \Omega_2$ が点 $x \in \Omega_1$ で\textbf{連続}であるとは，
  \[
    \forall \varepsilon > 0,\; \exists \delta > 0,\;
    d_1(x, y) < \delta \Longrightarrow d_2(f(x), f(y)) < \varepsilon
  \]
  が成り立つことをいう．$f$ が $\Omega_1$ の全点で連続であるとき，
  $f$ は\textbf{連続写像}であるという．
  同値な条件として，$\Omega_1$ の任意の収束列 $x_n \to x$ に対して
  $f(x_n) \to f(x)$ が成り立つことが挙げられる．
\end{definition}

\begin{proposition}{連続写像による閉集合の引き戻し}{sem-preimage-closed}
  距離空間 $(\Omega_1, d_1)$，$(\Omega_2, d_2)$ の間の連続写像 $f \colon \Omega_1 \to \Omega_2$ と
  閉集合 $F \subset \Omega_2$ に対して，逆像
  \[
    f^{-1}(F) \defeq \{ x \in \Omega_1 \mid f(x) \in F \}
  \]
  は $\Omega_1$ の閉集合である．
\end{proposition}
\begin{proof}
  $f^{-1}(F)$ に含まれる任意の収束列 $x_n \to x$（$x \in \Omega_1$）をとる．
  $f$ の連続性（定義~\ref{def:sem-continuous}の同値条件）より $f(x_n) \to f(x)$．
  仮定より $f(x_n) \in F$ であり，距離空間における閉集合は収束列の極限を含むから
  $f(x) \in F$，すなわち $x \in f^{-1}(F)$．
\end{proof}

\begin{definition}{有界集合}{sem-bounded}
  $\R^n$ の部分集合 $S \subset \R^n$ が\textbf{有界}であるとは，
  ある $M > 0$ が存在して
  \[
    \forall x \in S,\quad \|x\| \leq M
  \]
  が成り立つことをいう．ここで $\|\cdot\|$ はユークリッドノルムを表す．
\end{definition}

\begin{definition}{コンパクト集合}{sem-compact}
  距離空間 $(\Omega, d)$ の部分集合 $K \subset \Omega$ が
  \textbf{コンパクト}であるとは，
  $K$ における任意の点列 $(x_n)_{n \geq 1}$ が，
  $K$ の点に収束する部分列 $(x_{n_k})_{k \geq 1}$ を持つことをいう
  （列コンパクト性）．
\end{definition}

\begin{claim}{Bolzano-Weierstrass の定理}{sem-bolzano-weierstrass}
  $\R^n$ における任意の有界列 $(x_n)_{n \geq 1} \subset \R^n$ は，
  収束する部分列 $(x_{n_k})_{k \geq 1}$ を持つ．
\end{claim}

\begin{theorem}{Heine-Borel の定理}{sem-heine-borel}
  $\R^n$ の部分集合 $K \subset \R^n$ に対して，
  次の2条件は同値である：
  \begin{enumerate}
    \item[(i)] $K$ はコンパクト（定義~\ref{def:sem-compact}）である．
    \item[(ii)] $K$ は閉集合（定義~\ref{def:sem-topological-space}）かつ
      有界（定義~\ref{def:sem-bounded}）である．
  \end{enumerate}
\end{theorem}

\begin{proof}
  \textbf{(i) $\Rightarrow$ (ii)：コンパクト $\Rightarrow$ 有界かつ閉．}

  \textit{有界性．}
  $K$ が有界でないと仮定する．すると各 $n \in \N$ に対して
  $\|x_n\| > n$ をみたす $x_n \in K$ が存在する．
  $K$ のコンパクト性より $(x_n)$ はある収束部分列 $(x_{n_k})$ を持ち，
  その極限 $x^* \in K$ について $\|x_{n_k}\| \to \|x^*\|$．
  一方 $\|x_{n_k}\| > n_k \to \infty$ であり矛盾．
  ゆえに $K$ は有界．

  \textit{閉性．}
  $K$ への収束列 $(x_n)$ を任意にとり，極限を $x^*$ とする（$x^* \in \R^n$）．
  $K$ のコンパクト性より $(x_n)$ はある収束部分列 $(x_{n_k})$ を持ち，
  その極限 $y^* \in K$ が存在する．
  部分列も元の列の収束先と同じ極限を持つから $y^* = x^*$，
  すなわち $x^* \in K$．ゆえに $K$ は閉集合．

  \textbf{(ii) $\Rightarrow$ (i)：有界かつ閉 $\Rightarrow$ コンパクト．}

  $K$ の任意の点列 $(x_n)$ をとる．
  $K$ は有界ゆえ $(x_n)$ は $\R^n$ における有界列である．
  Bolzano-Weierstrass の定理（主張~\ref{clm:sem-bolzano-weierstrass}）より
  $(x_n)$ はある収束部分列 $(x_{n_k})$ を持ち，その極限を $x^* \in \R^n$ とする．
  $x_{n_k} \in K$ かつ $K$ は閉集合だから $x^* \in K$．
  ゆえに $K$ はコンパクト．
\end{proof}

\begin{claim}{部分列の極限がすべて等しいとき点列は収束する}{sem-subseq-limit-unique}
  コンパクト集合 $K \subset \R^n$ 内の点列 $(x_k)_{k \geq 1}$ と点 $x^\star \in K$ について，
  $(x_k)$ の収束する部分列がすべて $x^\star$ を極限に持つならば，
  $(x_k)$ 自身が $x^\star$ に収束する．
\end{claim}

\begin{proof}
  背理法による．$(x_k)$ が $x^\star$ に収束しないと仮定する．
  収束の定義の否定より，ある $\delta > 0$ と添字 $k_1 < k_2 < \cdots$ が存在して
  \[
    \|x_{k_j} - x^\star\| \geq \delta \qquad (\forall j)
  \]
  が成り立つ．$K$ はコンパクト（定義~\ref{def:sem-compact}）だから，
  この部分列 $(x_{k_j})_{j \geq 1}$ はさらに収束する部分列 $(x_{k_{j_\ell}})_{\ell \geq 1}$ を持ち，
  その極限を $y \in K$ とする．これは $(x_k)$ の収束部分列でもあるから，
  仮定よりその極限は $x^\star$ に等しく，$y = x^\star$．
  一方 $\|x_{k_{j_\ell}} - x^\star\| \geq \delta$ で $\ell \to \infty$ とすると
  $\|y - x^\star\| \geq \delta > 0$，すなわち $y \neq x^\star$ となり矛盾する．
\end{proof}

\begin{theorem}{Weierstrass の最大値の定理}{sem-weierstrass}
  $K \subset \R^n$ を空でないコンパクト集合，
  $f \colon K \to \R$ を連続関数（定義~\ref{def:sem-continuous}）とする．
  このとき $f$ は $K$ 上で最小値および最大値を達成する，すなわち
  \[
    \exists x^* \in K,\quad f(x^*) = \min_{x \in K} f(x).
  \]
\end{theorem}

\begin{proof}
  最小値について示す（最大値は $-f$ に適用すればよい）．

  $p^* \defeq \inf_{x \in K} f(x) \in [-\infty, +\infty)$ とおく．
  下限の定義より，各 $n \in \N$ に対して
  \[
    f(x_n) < p^* + \tfrac{1}{n}
  \]
  をみたす $x_n \in K$ が存在する（最小化列）．
  $K$ のコンパクト性より $(x_n)$ はある収束部分列 $(x_{n_k})$ を持ち，
  極限 $x^* \in K$ が存在する．
  $f$ の連続性（定義~\ref{def:sem-continuous}）より
  \[
    f(x^*) = \lim_{k \to \infty} f(x_{n_k}) = p^*.
  \]
  特に $p^* = f(x^*) \in \R$ であり，最小値は $x^*$ で達成される．
\end{proof}

\begin{proposition}{有限次元ノルム空間上の線形関数の連続性}{sem-finite-dim-linear-continuous}
  有限次元ノルム空間 $V$ 上の任意の線形関数 $f \colon V \to \R$ は連続である．
\end{proposition}

\begin{definition}{微分可能性と導関数}{sem-differentiable}
  開区間 $(a, b) \subset \R$ 上の関数 $f \colon (a, b) \to \R$ が
  点 $t \in (a, b)$ で\textbf{微分可能}であるとは，極限
  \[
    f'(t) \defeq \lim_{h \to 0} \frac{f(t + h) - f(t)}{h}
  \]
  が（有限な実数として）存在することをいい，この値 $f'(t)$ を
  $f$ の点 $t$ における\textbf{微分係数}という．
  $f$ が $(a, b)$ のすべての点で微分可能であるとき，
  $f$ は $(a, b)$ 上で\textbf{微分可能}であるといい，
  各点 $t$ に $f'(t)$ を対応させる関数 $f' \colon (a, b) \to \R$ を
  $f$ の\textbf{導関数}という．
\end{definition}

\begin{theorem}{平均値定理}{sem-mvt}
  $f \colon [a, b] \to \R$ が連続で $(a, b)$ 上で微分可能
  （定義~\ref{def:sem-differentiable}）であるとき，
  \[
    f(b) - f(a) = f'(c)(b - a)
  \]
  をみたす $c \in (a, b)$ が存在する．
\end{theorem}

\begin{remark}{有限集合は自動的に Polish}{sem-finite-auto-polish}
  距離空間 $(\Omega, d)$ の有限部分集合 $S = \{x_1, \ldots, x_n\} \subset \Omega$ は，
  誘導距離に関して自動的に Polish 空間（完備かつ可分，定義~\ref{def:sem-polish}）となる．

  \textbf{完備性．} $n = 1$ なら $S = \{x_1\}$ ゆえ任意の列は定値列 $y_k = x_1$ であり
  $x_1$ に収束する．$n \geq 2$ のときは，相異なる 2 点間の距離の最小値
  $\delta \defeq \min_{i \neq j} d(x_i, x_j) > 0$ をとる．任意の Cauchy 列
  $(y_k) \subset S$ に対し $\varepsilon \defeq \delta/2$ を取ると，ある $N \in \N$ 以降
  $d(y_k, y_l) < \delta/2 < \delta$ となるが，$S$ の相異なる 2 点の距離は $\delta$ 以上
  だからこれは $y_k = y_l$ を強制する．ゆえに $(y_k)$ は $N$ 以降一定の点に収束する．

  \textbf{可分性．} $S$ 自身が有限（特に可算）集合であり，自身の中で稠密だから可分．
\end{remark}
